% Section 7 (RQ3). Every number that comes from the corpus-wide ablation harvest is a macro defined in
% tables/rq3_macros.tex, generated together with the three table fragments input below, so a rerun of
% the harvest updates this section without editing it. So are the two sensitivity analyses (the layer
% correlation before the mapping rules, the cross-paper headline without baseline-only keys). Numbers
% that do not come from that harvest (the other cross-paper contrasts, the coded-set harvest, the
% own-arm test, the attribution study, the pilot) are literals and are fixed.
% Macro formats this text assumes (all used in text mode, never inside $...$):
%   counts            plain numbers, thousands separator where needed (\rqCorpusContrasts -> 867)
%   *Pct              number only, no percent sign; the text appends \% (\rqSignSharePct -> 91)
%   \rqSVPooled, \rqSVDiscounted, \rqCoverageRho, \rqCoverageP   signed or plain decimal, text-safe
%   \rqPreRuleRho, \rqPreRuleP, \rqXPOtherEffect, \rqXPOtherP     likewise
%   \rqSVCIz, \rqSVPI, \rqXPOtherCI  interval including its brackets, text-safe ([$+0.107$, $+0.187$])
%   \rqZeroDimList, \rqPIExcludesZeroDims   \texttt names joined with commas and "and"
%   \rqZeroLayerList  layer letters joined with "and"; the word "empty" if the set is empty
%                     (the sentence reads "the set of layers with no ablated dimension is \rqZeroLayerList")
%   \rqHarvestTierCoverage   a noun phrase (e.g. "28\% of the top-signal tier and none of the second")
%   \rqHarvestDate    a date
%   \rqSignSharePct   share of the \rqSignDenom contrasts on POOLED dimensions (coded-set plus
%                     corpus-wide) that favour the component

\section{Design choices versus outcomes}\label{sec:outcomes}

Research question~3 asks which harness design choices are associated with better task outcomes.
Three designs can address it, and they differ in what each can identify
(Figure~\ref{fig:rq3-triangulation}). A cross-paper association compares coded systems inside a shared
benchmark, split and base model; measurement error on the design variable attenuates it towards zero,
while the choice of reference systems can inflate it. A within-study meta-analysis pools the ablations
authors publish about their own components; selective reporting inflates it, so its estimates are
\emph{upper bounds}. Only a controlled, compute-matched ablation identifies a component's effect at
equal compute. Table~\ref{tab:rq3-bounds} sets the three side by side with the direction of each
bound. \citet{guo2026qa2task} run the closest prior
analysis; they compare harnesses by identity, where we compare them by coded dimension.

\begin{figure}[t]
  \centering
  \includegraphics[width=\linewidth]{rq3_triangulation}
  \caption{The three RQ3 designs, placed by what each can identify. Cross-paper association within comparable keys (36 keys) can show only co-occurrence; scaffold drift biases it towards zero and baseline selection away from it, so the $+0.88$ within-key sd for multi-agent topology is a lower bound only against drift (Section~\ref{subsec:comparability}). The within-study meta-analysis of authors' own ablations identifies a component effect but is inflated by unreported nulls, so its pooled effects are upper bounds; the corpus-wide row pools \rqTotalContrasts{} contrasts from \rqTotalPapers{} papers under the same design. Only a controlled compute-matched ablation identifies the effect at equal compute; our filed one failed its pilot band, so the one design that could settle RQ3 has not yet produced an estimate.}
  \label{fig:rq3-triangulation}
\end{figure}

% RQ3 designs and their bounds
% GENERATED by scripts/make_rq3_tables.py - do not edit; re-run the script instead.
% sources (mtime, UTC):
%   paper/tables/outcomes_summary.json  (2026-09-25T15:27:23Z)
%   data/analysis/ablation_pooled_corpus.csv  (2026-09-25T16:13:00Z)
%   data/analysis/ablation_credible_corpus.csv  (2026-09-25T16:13:00Z)
%   data/analysis/ablation_summary_corpus.json  (2026-09-25T16:13:00Z)
%   data/tier3/pilot/pilot_summary.json  (2026-09-25T06:14:15Z)
%   data/tier3/pilot_v2/pilot_v2_summary.json  (2026-09-25T15:56:13Z)
%   schema/dimensions.json  (2026-09-25T14:21:14Z)
% generated: 2026-09-25T17:21:54Z
\begin{table}
\centering
\caption{Three designs for RQ3 and the direction of their bounds. The within-study and registered designs are shown for self-verification (E1), the one component all three measure; the cross-paper design shows its headline contrast (multi-agent topology (C3), the one that survives both the key bootstrap and the permutation test) and, beneath it, the focal dimension's own cross-paper value. The cross-paper estimate is a standardised within-key difference ($d$, in within-key SDs); the within-study estimate is a relative change. They are not on one scale; what they share is the direction in which each can be wrong.}
\label{tab:rq3-bounds}
\scriptsize
\setlength{\tabcolsep}{3pt}
\begin{tabular}{@{}p{0.10\linewidth}p{0.16\linewidth}p{0.15\linewidth}p{0.12\linewidth}p{0.12\linewidth}p{0.09\linewidth}p{0.14\linewidth}@{}}
\toprule
Design & Identifies & $n$ & Headline estimate & Interval (95\%) & Bound & Status \\
\midrule
% BEGIN ROWS
\raggedright Cross-paper association & \raggedright association between a coded design choice and standardised score across systems sharing benchmark, split and base model & \raggedright 18 keys, 38 systems (of 36 keys, 53 systems) & \raggedright multi-agent topology (C3): $d = +0.877$ within-key SD & \raggedright \ensuremath{[+0.379, +1.313]} (bootstrap over keys) & \raggedright lower for scaffold drift; in 7 of its 18 keys baseline selection can inflate it & \raggedright complete; detected; exact permutation $p = 0.0013$ \tabularnewline
\raggedright  & \raggedright  & \raggedright 3 keys, 11 systems & \raggedright self-verification (E1), executable checks vs none; fragile (3 keys): $d = -0.734$ within-key SD & \raggedright \ensuremath{[-1.414, -0.258]} (bootstrap over keys) & \raggedright lower (attenuated towards zero by scaffold drift) & \raggedright complete; not detected; exact permutation $p = 0.306$; fragile \tabularnewline
\raggedright Within-study meta-analysis & \raggedright the component's contribution inside the host systems whose authors ablated it & \raggedright 213 papers, 680 contrasts (all dimensions: 494 papers, 2{,}466 contrasts) & \raggedright self-verification (E1): $\hat\mu = +0.157$ (relative change) & \raggedright $z$ \ensuremath{[+0.132, +0.183]}; HK \ensuremath{[+0.131, +0.183]} & \raggedright upper (selective reporting) & \raggedright interim (harvest has read 96.7\% of the top-signal tier and none of the second); survives the bias discount ($z$ and HK) \tabularnewline
\raggedright Registered compute-matched ablation & \raggedright causal effect of the component at a matched model-call budget, on one suite and model & \raggedright pilot: 9 suite--model cells, arm A; revision: 48 runs, arms A and B; 93 instances needed at $\rho = 0.7$, 22 available & \raggedright not yet estimated & \raggedright --- & \raggedright unbiased (for that suite and model) & \raggedright filed, amendment pending; pilot failed its band; not runnable on available suites \tabularnewline
% END ROWS
\bottomrule
\end{tabular}
\end{table}


\subsection{Fewer than one coded system in twenty can be compared to a peer}\label{subsec:comparability}

The cross-paper design is nearly empty. Of the 1,256 coded systems the outcome analysis read, 646
report a benchmark score, but only 53 (4.2\%) share a canonical benchmark, split and base model with
another coded system, across 36 comparable keys. The registered mixed-effects regression was therefore
not fitted (amendment~11): its 38 design dimensions against 36 key clusters leave the coefficients
without support, and we report the non-fit rather than substitute an unregistered model.

Two of five contrasts, specified in amendment~11 after the comparable set was assembled, clear both of
their tests, and neither is free of selection.
Multi-agent topology is associated with $+0.88$ within-key standard deviations, 95\% bootstrap CI
$[+0.38, +1.31]$ over keys, on 18 keys and 38 systems, within-key permutation $p=0.0013$. It is an
association, not an effect, and two biases act on it in opposite directions. Scores do not name the
harness commit that produced them, while we code a pinned commit; that scaffold drift is measurement
error on the design variable and attenuates the estimate. Selection
pushes the other way: in \rqXPBaselineKeys{} of the 18 keys every reference system is a recurring
baseline harness, and on the other \rqXPOtherKeys{} keys the association is \rqXPOtherEffect{}
\rqXPOtherCI{} (permutation $p$ of \rqXPOtherP), so it is a lower bound only against drift, and its
detection rests on the baseline keys.

The other contrasts are reported with what they could have detected. \texttt{explicit\_plan} is a null,
$+0.36$ $[-0.30, +0.97]$, with a minimum detectable effect (MDE) of 1.05 within-key standard
deviations. \texttt{executable\_verification}, on 3 keys, has a bootstrap interval that excludes zero,
$-0.73$ $[-1.41, -0.26]$, but an exact permutation test that does not ($p=0.306$; MDE 1.87), so it is
read as a null. \texttt{any\_retry} varies inside only 2 keys, whose sign-symmetric label assignments
floor its two-sided $p$ at 0.500, so it is not estimable. \texttt{any\_compaction} clears both tests at
its design floor ($p=0.042$), but all 4 of its keys use a recurring baseline as reference, so it is
reported as confounded.

\emph{Takeaway.} The cross-paper literature supports one association, detected only with the keys
that compare against baselines, and cannot identify design effects, a measured property of how the
field reports results.

\subsection{Authors' own ablations barely separate design dimensions}\label{subsec:ablations}

The within-study design pools authors' published ablations of their own systems. Each contrast is one
codebase at one commit with one base model, differing in one component, and its effect is the relative
change $(\text{with}-\text{without})/\text{with}$ in the paper's own metric. Beyond the
\rqCodedContrasts{} coded-set contrasts from \rqCodedPapers{} papers, the completed extraction stage
read the full text of \rqCorpusPapersExtracted{} included papers, \rqHarvestTierCoverage{} of a
prefilter that ranks papers by ablation content, and added \rqCorpusContrasts{} pool-eligible contrasts from \rqCorpusPapersWithContrast{}
papers. \rqPoolableDims{} dimensions now reach the three papers needed to pool, against 8 on the
coded set alone.

Rows enter the pool only under fixed rules. A row is kept if its two scores appear together and
verbatim, share a scale and metric, vary the harness rather than the model, and map to a design
dimension above a confidence floor. Seven mapping rules then correct attributions the text
contradicts: a removed sandbox that takes code-execution feedback with it is self-verification, not
isolation, and a row whose scores run together, report a delta, or carry a label that contradicts its
direction is dropped. Every firing is logged in the replication package.

% RQ3 pooled ablations
% GENERATED by scripts/make_rq3_tables.py - do not edit; re-run the script instead.
% sources (mtime, UTC):
%   data/analysis/ablation_pooled_corpus.csv  (2026-09-25T16:13:00Z)
%   data/analysis/ablation_credible_corpus.csv  (2026-09-25T16:13:00Z)
%   data/analysis/ablation_summary_corpus.json  (2026-09-25T16:13:00Z)
%   schema/dimensions.json  (2026-09-25T14:21:14Z)
% generated: 2026-09-25T16:13:01Z
\begin{table}
\centering
\caption{Within-study meta-analysis of published ablations, coded-set and corpus-wide harvests combined (harvest snapshot of 25 September 2026; the harvest is incomplete). One row per design dimension with contrasts from at least 3 papers (18 dimensions; 6 more have contrasts from fewer papers and are not pooled). Effect: relative change $(\mathrm{with}-\mathrm{without})/\mathrm{with}$ on the paper's own metric, $+$ = the component helped. $\hat\mu$: DerSimonian--Laird random-effects mean over paper-level effects; 95\% CI by the normal ($z$) and Hartung--Knapp (HK) methods; 95\% PI: prediction interval for a new host system, \textbf{bold} where it excludes zero. Discounted $\hat\mu$ applies the dimension's own publication-bias discount. \emph{Surv.}: the discounted lower CI bound still exceeds 0.02 and trim-and-fill keeps the sign, judged on the $z$ / HK bound. Papers: total (coded-set + corpus-wide). Every estimate is an \emph{upper} bound on the component's contribution: authors report the ablations that favour their component.}
\label{tab:ablations}
\footnotesize
\setlength{\tabcolsep}{2.5pt}
\begin{tabular}{@{}p{0.2\linewidth}rrrcccrrc@{}}
\toprule
 & & & \multicolumn{7}{c}{relative change, $+$ = component helped} \\
\cmidrule(l){4-10}
Dimension & Papers & Contr. & $\hat\mu$ & 95\% CI ($z$) & 95\% CI (HK) & 95\% PI & $I^2$ (\%) & Disc.\ $\hat\mu$ & Surv.\ $z$/HK \\
\midrule
% BEGIN ROWS
\raggedright Self-verification (E1) & 213 (14+201) & 680 & \ensuremath{+0.157} & \ensuremath{[+0.132,+0.183]} & \ensuremath{[+0.131,+0.183]} & \ensuremath{[-0.172,+0.487]} & 91.7 & \ensuremath{+0.082} & yes/yes \tabularnewline
\raggedright Long-term memory (D2) & 133 (4+131) & 455 & \ensuremath{+0.151} & \ensuremath{[+0.125,+0.176]} & \ensuremath{[+0.123,+0.179]} & \ensuremath{[-0.095,+0.397]} & 77.3 & \ensuremath{+0.080} & yes/yes \tabularnewline
\raggedright Planning granularity (C2) & 86 (4+83) & 245 & \ensuremath{+0.161} & \ensuremath{[+0.126,+0.197]} & \ensuremath{[+0.127,+0.196]} & \ensuremath{[-0.128,+0.451]} & 81.6 & \ensuremath{+0.085} & yes/yes \tabularnewline
\raggedright Repo/environment context strategy (A2) & 85 (5+80) & 253 & \ensuremath{+0.177} & \ensuremath{[+0.129,+0.226]} & \ensuremath{[+0.128,+0.227]} & \ensuremath{[-0.241,+0.596]} & 90.4 & \ensuremath{+0.090} & yes/yes \tabularnewline
\raggedright Multi-agent topology (C3) & 79 (6+74) & 272 & \ensuremath{+0.174} & \ensuremath{[+0.130,+0.218]} & \ensuremath{[+0.129,+0.219]} & \ensuremath{[-0.181,+0.530]} & 88.6 & \ensuremath{+0.090} & yes/yes \tabularnewline
\raggedright Loop primitive(s) (C1) & 52 (1+51) & 156 & \ensuremath{+0.192} & \ensuremath{[+0.139,+0.245]} & \ensuremath{[+0.148,+0.236]} & \ensuremath{[-0.163,+0.547]} & 87.6 & \ensuremath{+0.099} & yes/yes \tabularnewline
\raggedright Short-term state (D1) & 30 (3+28) & 59 & \ensuremath{+0.206} & \ensuremath{[+0.108,+0.304]} & \ensuremath{[+0.098,+0.314]} & \ensuremath{[-0.330,+0.742]} & 91.6 & \ensuremath{+0.104} & yes/yes \tabularnewline
\raggedright Tool count at pinned version (B2) & 27 (4+23) & 95 & \ensuremath{+0.173} & \ensuremath{[+0.017,+0.328]} & \ensuremath{[+0.079,+0.267]} & \ensuremath{[-0.675,+1.021]} & 97.9 & \ensuremath{+0.087} & no/yes \tabularnewline
\raggedright Context compaction (A3) & 19 (4+15) & 43 & \ensuremath{+0.085} & \ensuremath{[+0.048,+0.122]} & \ensuremath{[+0.046,+0.124]} & \textbf{\ensuremath{[+0.029,+0.141]}} & 5.2 & \ensuremath{+0.052} & yes/yes \tabularnewline
\raggedright Delegation mechanism (C4) & 14 (1+13) & 42 & \ensuremath{+0.087} & \ensuremath{[-0.024,+0.197]} & \ensuremath{[-0.047,+0.221]} & \ensuremath{[-0.355,+0.528]} & 89.4 & \ensuremath{+0.002} & no/no \tabularnewline
\raggedright Observation formatting (A4) & 14 (1+13) & 37 & \ensuremath{+0.169} & \ensuremath{[+0.061,+0.277]} & \ensuremath{[+0.081,+0.257]} & \ensuremath{[-0.258,+0.596]} & 90.5 & \ensuremath{+0.087} & yes/yes \tabularnewline
\raggedright System-prompt style (A1) & 12 (1+11) & 31 & \ensuremath{+0.138} & \ensuremath{[+0.048,+0.227]} & \ensuremath{[+0.004,+0.272]} & \ensuremath{[-0.153,+0.428]} & 71.6 & \ensuremath{+0.067} & yes/no \tabularnewline
\raggedright Retry policy (E2) & 12 (1+11) & 23 & \ensuremath{+0.084} & \ensuremath{[+0.041,+0.128]} & \ensuremath{[+0.035,+0.133]} & \ensuremath{[-0.032,+0.201]} & 47.1 & \ensuremath{+0.046} & yes/no \tabularnewline
\raggedright Rollback / undo (E3) & 8 (1+7) & 21 & \ensuremath{+0.066} & \ensuremath{[-0.004,+0.137]} & \ensuremath{[-0.019,+0.151]} & \ensuremath{[-0.116,+0.249]} & 45.1 & \ensuremath{+0.037} & no/no \tabularnewline
\raggedright Guardrails (H4) & 7 (0+7) & 8 & \ensuremath{+0.107} & \ensuremath{[+0.069,+0.145]} & \ensuremath{[+0.067,+0.147]} & \textbf{\ensuremath{[+0.058,+0.157]}} & 0.0 & \ensuremath{+0.085} & yes/yes \tabularnewline
\raggedright Edit primitive (B3) & 4 (0+4) & 11 & \ensuremath{+0.169} & \ensuremath{[+0.042,+0.295]} & \ensuremath{[-0.046,+0.384]} & \ensuremath{[-0.261,+0.599]} & 34.8 & \ensuremath{+0.084} & yes/no \tabularnewline
\raggedright Termination condition (F1) & 4 (0+4) & 9 & \ensuremath{-0.007} & \ensuremath{[-0.025,+0.010]} & \ensuremath{[-0.018,+0.003]} & \ensuremath{[-0.045,+0.031]} & 0.0 & \ensuremath{-0.007} & no/no \tabularnewline
\raggedright Tracing (H1) & 3 (0+3) & 6 & \ensuremath{+0.052} & \ensuremath{[-0.050,+0.154]} & \ensuremath{[-0.151,+0.255]} & \ensuremath{[-1.014,+1.118]} & 46.1 & \ensuremath{+0.002} & no/no \tabularnewline
% END ROWS
\bottomrule
\end{tabular}
\end{table}


Table~\ref{tab:ablations} gives the pooled result, and its central fact is convergence. Most
dimensions show a positive relative gain of similar size: surviving gains span
\rqEffectRangeSurviving{}, their middle half \rqEffectIQRSurviving, and most intervals overlap.
Context compaction has the lowest surviving gain and the least heterogeneity of the well-populated
rows; its $z$ interval lies below those of multi-agent topology and loop primitives but overlaps the
other four best-populated dimensions. Short-term state, as well populated, pools near zero. The bias
signature repeats: \rqSignSharePct\% of the \rqSignDenom{} contrasts on pooled dimensions favour the
authors' component, heterogeneity is mostly high, and all \rqPIIncludesZeroCount{} prediction
intervals include zero.

Gains of one size from components this different fit one reporting filter better than coincidence:
an ablation is written up to show that a component helps. Misassigned labels would also pull rows
together, so convergence shows how little these ablations separate dimensions, not why.

The pre-stated discount rule fails to discriminate at this scale. The rule
shrinks each pooled effect by the larger of its own trim-and-fill and one-null-per-contrast shrinkage
\citep{duval2000trimfill,copas2000funnel}, and admits a dimension whose discounted confidence-interval
lower bound still clears 2\% relative. Of \rqPoolableDims{} dimensions, \rqDiscountSurviveZ{} pass, and
\rqDiscountSurviveHK{} pass with a Hartung--Knapp interval in place of the $z$ interval. The cause is
structural. For most positive effects the discount sits near one half (median \rqMedianDiscountPct\%), so
the rule reduces to asking whether an interval clears zero, and more contrasts make that easier
whatever produced the effect. On the coded set, with 3 to 14 papers per dimension, the same rule
admitted two dimensions; its selectivity there was a property of small samples.

The prediction-interval criterion \citep{higgins2009predictioninterval} is the secondary we can defend,
and it admits no dimension. A prediction interval clears zero only where between-paper heterogeneity is small, so the criterion asks
whether the next published ablation would also favour the component.

Self-verification's estimate is stable and no longer distinguished. On the coded set it pooled to at
most $+0.139$ on 14 papers ($+0.069$ after the discount); corpus-wide it pools to \rqSVPooled, 95\% CI
\rqSVCIz, on \rqSVPapers{} papers, with a prediction interval of \rqSVPI{} that includes zero, and
discounts to \rqSVDiscounted{} under one unreported null per contrast. A filtered literature produces
that stability as readily as a real effect does. The defensible sentence is that self-verification is worth at most about
\rqSVDiscounted{} relative in published ablations, and a sentence of the same form can be written for
most of the table.

Every estimate in Table~\ref{tab:ablations} is an upper bound. The within-study design absorbs four of
the five benchmark caveats we document, because both arms share one codebase, commit and team; what
survives is self-reporting, concentrated: each contrast is an author reporting on a component they
built.

\emph{Takeaway.} Authors' own ablations bound most components from above and rank almost none:
whatever is ablated appears to help by about the same amount, which fits selective reporting better
than equal real effects.

\subsection{Publication bias, measured on the sign distribution}\label{subsec:bias}

Selective reporting is measured here rather than assumed, and the measurement is the sign distribution.
Of the \rqSignDenom{} coded-set and corpus-wide contrasts on pooled dimensions, \rqSignSharePct\% favour
the authors' own component. The share does not rest on
our orientation of the metrics: \rqPolarityAssumedN{} contrasts take higher-is-better by assumption
because their metric names no direction, yet \rqPolarityStatedFavourPct\% of the
\rqPolarityStatedN{} with a stated direction favour the component (\rqPolarityAssumedFavourPct\% of the
rest). No rule reads an effect's sign, but the rule that drops a row whose label contradicts
its direction removed \rqDCLDropped{} rows (\rqDCLAgainst{} unfavourable as extracted); kept, they
would give \rqSignShareWithDCLPct\%. A component that helped only sometimes would produce more
dissent than that, and this distribution is the empirical warrant for reading
Table~\ref{tab:ablations} as upper bounds.

The within-paper baseline tables corroborate the mechanism; they are not a second measurement of it.
Over all 75 blocks containing the authors' own arm, from 37 papers, the own arm's advantage over chance
is not established (paper-clustered $p=0.21$). It survives only in tables of five or more arms, 14 of
17 against 6.95 expected ($p=4.4\times10^{-4}$), and those tables hold 70 of the 89 own-ablation arms:
large arm tables here are largely the paper's own ablation tables.

Egger's test \citep{egger1997bias} is reported but not leaned on: its input is partly our own
construction. Papers publish no standard errors, so a contrast's variance is reconstructed binomially
from its arm scores where the metric is a proportion, and set by an assumed 10\% relative standard
error for the other \rqVarAssumedN{} contrasts on pooled dimensions. Under that reconstruction,
at a given with-component score, the variance of the relative effect falls as the effect grows, so
random-effects weights \citep{dersimonian1986meta} favour the largest effects, pushing pooled estimates
upward as the upper bound declares, and the funnel tilts against small-study asymmetry, biasing
Egger's test towards non-detection.

\emph{Takeaway.} Author selection of what to publish is present at the magnitude the sign distribution
measures, and it acts on which contrasts reach print.

\subsection{Where the field ablates}\label{subsec:coverage}

Ablation density falls as a layer's silence rises. Table~\ref{tab:coverage} sets the harvested contrasts
against each layer's weighted \texttt{not\_reported} rate from the landscape analysis. Across the eight
substantive layers, the rank correlation between
silence and ablation density is Spearman's $\rho$ of \rqCoverageRho, exact one-sided permutation $p$ of
\rqCoverageP. On the coded set, layers F, G and H (budget and termination, sandbox and environment,
observability and governance) carried no ablation at all. The corpus-wide harvest reaches four of their
dimensions, through 1 to 21 papers each: termination condition (F1) pools to a precise null,
$-0.003$ $[-0.019, +0.014]$ with $I^2 = 0$, and guardrails (H4) to $-0.100$ with an interval spanning
zero, but \rqZeroDimCount{} dimensions still carry no published ablation (\rqZeroDimList), and the set
of layers with no ablated dimension is \rqZeroLayerList{}.

The mapping rules produce that empty layer but not the gradient. Of the \rqPreRuleGContrasts{}
contrasts the harvest first attributed to the sandbox layer, the rules reassigned
\rqPreRuleGRemapped{}, removals of code-execution feedback or file-system tools, to the dimensions they
take away (\rqPreRuleGRemapList), and dropped \rqPreRuleGDropped{}: \rqPreRuleGDroppedUnread{} with
unreadable scores and \rqPreRuleGDroppedUnnamed{} naming no isolation boundary or authorization. With
all seven rules undone $\rho$ is unchanged, \rqPreRuleRho{} (exact $p$ of \rqPreRuleP).

% RQ3 ablation coverage
% GENERATED by scripts/make_rq3_tables.py - do not edit; re-run the script instead.
% sources (mtime, UTC):
%   data/analysis/ablation_coverage_by_layer.csv  (2026-09-25T16:01:11Z)
%   data/analysis/ablation_coverage_summary.json  (2026-09-25T16:01:11Z)
%   data/analysis/ablation_coverage_by_layer_corpus.csv  (2026-09-28T13:53:57Z)
%   data/analysis/ablation_coverage_summary_corpus.json  (2026-09-28T13:53:57Z)
%   schema/dimensions.json  (2026-09-25T14:21:14Z)
% generated: 2026-09-28T13:53:57Z
\begin{table}
\centering
\caption{Ablation coverage by layer against layer silence (weighted share of systems whose dimension is \emph{not reported}). \emph{Coded-set}: contrasts harvested from the coded systems' outcome extraction; \emph{corpus-wide}: coded-set plus the full-text harvest, de-duplicated. Dimension and density columns are corpus-wide. $^\dagger$Layer M (metadata) cannot be ablated and is excluded from $\rho$. Still without a single published ablation: 7 of 31 ablatable dimensions (cost controls (F2), evaluation hooks (H3), execution isolation (G1), filesystem access (G2), network policy (G3), permission model (G4) and replayability (H2)). Absence of evidence about the literature, not evidence that those components do not matter.}
\label{tab:coverage}
\footnotesize
\begin{tabular}{@{}lrrrrrr@{}}
\toprule
 & Weighted & & Dims $\geq$1 & Contrasts & Contrasts & Contrasts \\
Layer & silence & Dims & contrast & (coded-set) & (corpus-wide) & per dim \\
\midrule
% BEGIN ROWS
A Context assembly & 40.2\% & 4 & 4 & 13 & 924 & 231.0 \\
B Tool interface & 71.4\% & 5 & 5 & 5 & 272 & 54.4 \\
C Control loop & 24.9\% & 5 & 5 & 26 & 1{,}651 & 330.2 \\
D Memory and state & 63.3\% & 3 & 3 & 8 & 1{,}263 & 421.0 \\
E Verification and repair & 60.9\% & 3 & 3 & 23 & 1{,}817 & 605.7 \\
F Budget and termination & 65.1\% & 3 & 2 & 0 & 29 & 9.7 \\
G Sandbox and environment & 81.6\% & 4 & 0 & 0 & 0 & 0.0 \\
H Observability and governance & 72.8\% & 4 & 2 & 0 & 47 & 11.8 \\
M Meta$^\dagger$ & 34.6\% & 7 & 0 & 0 & 0 & 0.0 \\
% END ROWS
\midrule
\multicolumn{7}{@{}l}{Spearman $\rho$, silence vs.\ corpus-wide density (layers A--H): $\rho = -0.738$, exact one-sided $p = 0.023$} \\
\multicolumn{7}{@{}l}{(40{,}320 permutations); coded-set density only: $\rho = -0.854$, $p = 0.0065$} \\
\bottomrule
\end{tabular}
\end{table}


This is an absence of evidence and is read as one: no paper in the harvest published a number for
removing these components, which says nothing about whether they matter. The layers that govern safe
deployment are the layers the literature has measured least.

\emph{Takeaway.} The field ablates the layers that describe what a harness does and seldom the layers
that bound it.

\subsection{Papers compare their system against its own configurations}\label{subsec:attribution}

The within-paper baseline table should have been RQ3's strongest design, and it fails for a reason
that is itself a finding. A block, one paper by benchmark by split by metric by model, holds
team, protocol and base model fixed. Yet 784 of the 819 comparator arms in the 324 estimable blocks
(95.7\%) cannot be tied to any system in the 6,504-system census, and the only 7 blocks with two or
more attributed comparators come from one paper, so no contrast in this design is replicated across
papers.

Name recall is not the explanation. A vetted alias table,
accepting a mapping only when the candidate's name is corroborated in that paper's own full text,
raised attributed comparator arms from 28 to 35 and left the number of blocks spanning more than one
paper unchanged. Of the 172 comparator labels the first reading judged to name a harness, 149 were
the reporting paper's own system under an abbreviation, such as \texttt{OH CodeActAgent v1.5} in the
OpenHands paper \citep{wang2024openhands}.

\emph{Takeaway.} Harness-to-harness comparison is largely not attempted; baseline
tables hold the paper's configurations and bare model names, a gap
\citet{samemodel2026harness} and \citet{scaffoldeffect2026} document from the model's side.

\subsection{Our own controlled ablation: filed, piloted, not yet run}\label{subsec:tier3}

Both published designs are author-reported, so before collecting any data we filed a controlled
ablation of our own as a protocol amendment, pending acceptance (Figure~\ref{fig:tier3-design}).
It targets self-verification, chosen when the coded-set harvest singled it out. A verification step
issues extra model calls, so every published contrast asks whether verifying helps, confounded with
compute, and never whether verifying and repairing beats unguided retrying at equal compute. Arm~C
therefore receives exactly the calls arm~B used, matched per instance rather than on average. No paper
among the 42 coded-set papers that publish a component ablation runs such a control, a distinction
that work on feedback compute makes salient \citep{efc2026scaling}.

\begin{figure}[t]
  \centering
  \includegraphics[width=\linewidth]{tier3_design}
  \caption{The filed compute-matched ablation, drawn for one instance. Arm~A makes a single call. Arm~B attempts, runs checks it wrote itself and repairs until they pass or $k$ calls are spent, recording the calls it used, $c_i$. Arm~C then receives exactly $c_i$ independent, unguided attempts on the same instance. The primary contrast, $B - C$ paired by instance, separates what verification does with its calls from how many calls it spends, which is the confound in every published self-verification ablation. Inset: the arm-A pilot against the pre-stated 25--70\% band. One of nine suite--model cells lies inside the band on its point estimate and none on its 95\% interval, so under the pre-stated rule no arm comparison is published.}
  \label{fig:tier3-design}
\end{figure}

The pilot failed its pre-stated band, so no arm comparison is published. The suite was to be chosen
on a disjoint pilot pool by arm-A accuracy in the 25--70\% band, read on the mean fraction of hidden
checks passed, and eight of the nine suite--model cells sit above the ceiling
(Figure~\ref{fig:tier3-design}, inset). The one inside, the Instruct presentation of BigCodeBench-Hard
with the smaller of the two generators, fails twice over: its interval
crosses the ceiling and cannot be resolved on a 23-instance pilot pool, and its confirmatory pool falls
far short of the 93 instances the analysis plan requires (Table~\ref{tab:rq3-bounds}). The full
Instruct split, which has the instances, sits entirely above the ceiling. Four of the nine cells are in
band on the binary secondary measure; we did not substitute it for the primary one.

One single-arm, descriptive pilot measurement is reportable. Across 56 arm-B diagnostic runs, a harness
that writes, runs and repairs against its own checks accepted its first attempt 50 times, at a mean of
1.16 calls, and three of those 50 scored 0.00 against the hidden tests. Model-written verification
mostly does not fire, and when it passes it can be blind to total failure, which bears on the mechanism
the published contrasts credit without contradicting their upper bound.

Where arm~B declines to spend its budget, arm~C is matched to a single call, and on those instances the
paired contrast carries no information about verification. A revision of the design, filed as a dated
amendment, removes that option: arm~B runs four verify-and-repair rounds whether or not its own checks
pass, and the estimand is conditioned on the stratum in which the treatment fires, defined by arm~A
rather than by arm~B's own outcome so that the conditioning cannot bias the contrast against~B. Its
pilot, on arms A and~B with the smaller generator, shows the redefinition works: the forced rounds
changed the submitted code in 17 of 48 runs, while the original arm accepted its first attempt in 41 of
48, three of them scoring 0.00 on the hidden tests.

The revision also exposed a defect that leaves the verdict unchanged. Pilot and confirmatory pools had
been separated by suite name rather than by task, so 12 of the in-band cell's 34 confirmatory tasks had
already been piloted under another presentation. Separated by task, 22 clean instances remain against
the 93 required; assuming the two arms' instance means correlate at 0.7, the smallest effect 22
instances can detect is 0.149 on the continuous score, over twice the planned 7 points. The design is
sound and still unrunnable on the suites available.
% TIER3V2-HOOK (filled 2026-09-25 from the amendment 12a pilot)

\emph{Takeaway.} The compute-matched question that would turn Section~\ref{subsec:ablations}'s upper
bounds into an unconfounded estimate is filed and unanswered.

% TODO: the protocol amendment for the compute-matched ablation is filed but not yet accepted at the
% registry. When it is, give it its number and date in the first paragraph of this subsection, in the
% amendments table of the supplement, and in the PRISMA checklist.
